\section{Kolmogorov Complexity}
\label{sec:kolmogorov}

\subsection{The Intuition}

Consider two strings of a thousand bits. String A alternates 0 and 1
indefinitely; String B is a typical output of a fair coin. Both are a
thousand bits long. String A can be described in a few dozen characters:
``alternating 0 and 1, one thousand times.'' String B, if genuinely random,
cannot be described more briefly than writing it out in full.

Kolmogorov complexity makes this precise. The complexity $\K(x)$ of a string
$x$ is the length of the shortest computer program that outputs $x$ and halts.
String A has very low $\K$; String B has $\K$ close to its own length.

Three examples locate the concept concretely. The digits of $\pi$ look random
but have very low $\K$: any standard algorithm for computing $\pi$ provides a
program of a few hundred bytes that generates arbitrarily many digits.
Newton's law of universal gravitation $F = Gm_1 m_2 / r^2$ is a program of a
few dozen bytes that predicts the trajectories of every planet, moon, and comet
we have observed, to many decimal places. A random genome, by contrast, would
have $\K$ close to its raw length---but the actual human genome is substantially
more compressible, because evolution is a compressive process.

\subsection{The Formal Definition}

Fix a universal Turing machine $U$. The Kolmogorov complexity of a finite
binary string $x$ with respect to $U$ is:
\[
  \K_U(x) = \min\{ |p| : U(p) = x \}
\]
where $|p|$ is the length of program $p$ in bits. In practice, the
prefix-free variant (where no valid program is a prefix of another) is
standard; it has cleaner information-theoretic properties.

The key theorems:

\begin{theorem}[Invariance~\cite{kolmogorov1965}]
For any two universal Turing machines $U$ and $V$:
\[
  |\K_U(x) - \K_V(x)| \leq c_{U,V}
\]
where $c_{U,V}$ depends only on the machines, not on $x$. For large strings,
this constant is negligible relative to $\K(x)$.
\end{theorem}

The Invariance Theorem is what makes $\K(x)$ objective: the choice of machine
changes the value by at most a fixed overhead, independent of the string being
measured. For any two computable descriptions, the complexity is the same up
to a translator constant.

\begin{proposition}
For any length $n$, at most $2^{n-c}$ strings of length $n$ have
$\K(x) < n - c$. Most strings are incompressible.
\end{proposition}

This is a counting argument: there are only $2^{n-c}$ programs shorter than
$n-c$ bits. When we encounter a compressible string in nature---a genome, a
physical trajectory, a spectrum---that compressibility is evidence of an
underlying generative process. Structure is rare.

\subsection{Uncomputability: A Feature, Not a Bug}

$\K(x)$ is not computable. No program, given $x$ as input, reliably outputs
$\K(x)$ and halts. The proof is a formalization of Berry's paradox: supposing
$\K$ were computable, one could construct a short program that outputs a string
with high complexity---a contradiction.

This is sometimes presented as a defect. It is not. Kolmogorov complexity is
the \emph{right} measure of structural complexity. Its uncomputability is the
signature of universality: $\K$ captures structure defined by all possible
programs, not any particular one. If $\K$ were computable, it would be
reducible to some specific algorithm, and that algorithm would make implicit
choices about what structure counts.

The scientist's relationship to $\K$ is the same as the mathematician's
relationship to real numbers: most are uncomputable, yet the reals are the
right framework. Scientists approximate $\K$ from above. Every proposed
theory is an upper bound on $\K(D)$. The search for shorter descriptions is
the search for a tighter bound. The fact that the true minimum is uncomputable
is precisely why the search never ends.

\subsection{MDL as the Computable Approximation}

The Minimum Description Length principle~\cite{rissanen1978} replaces $\K$
with a description length under a fixed coding scheme. For a parametric model
$H$ with parameters $\theta$:
\[
  L(H) + L(D \mid H) = L(\theta) - \log P(D \mid \theta)
\]
This is computable for any fixed model class and coding scheme. MDL is
model-selection consistent: as data grows, it converges to the true model
when the true model is in the hypothesis class. It reduces overfitting by
penalizing complexity directly, without ad hoc regularization.

The relationship to Bayes is exact. Under a prior $P(H) \propto 2^{-L(H)}$:
\[
  \log P(H \mid D) = -L(H) + \log P(D \mid H) + \text{const} = -\MDL(H,D) + \text{const}
\]
Maximizing the Bayesian posterior is equivalent to minimizing MDL, for the
right choice of prior and coding scheme.

\subsection{Connections to Information Theory}

The expected Kolmogorov complexity of a draw from a computable distribution
$P$ equals its Shannon entropy $H(P)$ up to lower-order terms:
\[
  \mathbb{E}[K(X)] \approx H(X).
\]
Shannon entropy is Kolmogorov complexity averaged over a distribution.
Kolmogorov complexity is the pointwise version: it applies to individual
strings, without requiring a distribution. When you have a single string and
want to know how much structure it has---without committing to a probability
model---Kolmogorov complexity is the right tool.

Lossless compression (gzip, bzip2, zstd) provides empirical upper bounds on
$\K$. A string that compresses well has low $\K$; one that does not compress
has high $\K$. This means compression ratio is a practical proxy for
structural complexity, usable without specifying any explicit model.